% p011_e 的电脑重画（等量异种电荷的电场线；左负右正，中垂线虚线标 φ=0，红线为 C 点两电荷场强方向与夹角 θ）
\begin{tikzpicture}[line width=0.7pt, >={Stealth[length=2.2mm]},
  fl/.style={postaction={decorate}, decoration={markings, mark=at position #1 with {\arrow{>}}}}]
  % 中垂线（虚线）
  \draw[dashed] (0,-2.25) -- (0,1.9);
  % 两电荷之间的三条电场线（箭头向左）
  \draw[fl=0.5] (1.3,0) -- (-1.3,0);
  \draw[fl=0.5] (1.36,0.15) .. controls (0.95,0.65) and (0.5,0.83) .. (0,0.83)
                            .. controls (-0.5,0.83) and (-0.95,0.65) .. (-1.36,0.15);
  \draw[fl=0.5] (1.36,-0.15) .. controls (0.95,-0.65) and (0.5,-0.83) .. (0,-0.83)
                             .. controls (-0.5,-0.83) and (-0.95,-0.65) .. (-1.36,-0.15);
  % 右侧正电荷发出的电场线（箭头向外）
  \draw[fl=0.55] (1.52,0.2)   .. controls (1.75,0.8) and (1.65,1.25) .. (1.33,1.5);
  \draw[fl=0.55] (1.64,0.14)  .. controls (2.2,0.5)   and (2.6,0.9)   .. (2.82,1.65);
  \draw[fl=0.62] (1.68,0.08)  .. controls (2.6,0.35)  and (3.1,0.8)   .. (3.62,1.55);
  \draw[fl=0.78] (1.7,0) -- (4.08,0);
  \draw[fl=0.6]  (1.69,-0.07) .. controls (2.6,-0.25) and (3.1,-0.6)  .. (3.2,-1.9);
  \draw[fl=0.45] (1.6,-0.17)  .. controls (2.1,-0.6)  and (2.3,-1.0)  .. (2.27,-2.2);
  \draw[fl=0.6]  (1.48,-0.2)  .. controls (1.6,-0.8)  and (1.3,-1.3)  .. (0.85,-1.83);
  % 左侧负电荷接收的电场线（箭头指向电荷）
  \draw[fl=0.5]  (-3.06,0) -- (-1.7,0);
  \draw[fl=0.7]  (-2.26,1.6)  .. controls (-2.45,0.9)  and (-2.2,0.4)   .. (-1.69,0.07);
  \draw[fl=0.6]  (-0.86,1.62) .. controls (-1.2,1.15)  and (-1.5,0.9)   .. (-1.55,0.2);
  \draw[fl=0.65] (-2.7,-1.76) .. controls (-2.4,-1.1)   and (-2.1,-0.6)  .. (-1.65,-0.15);
  \draw[fl=0.5]  (-1.68,-2.05) .. controls (-1.7,-1.3)  and (-1.55,-0.8) .. (-1.52,-0.2);
  % 两个电荷
  \draw[line width=1.4pt, fill=white] (-1.5,0) circle (0.2);
  \draw[line width=1.4pt] (-1.7,0) -- (-1.3,0);
  \draw[fill=white] (1.5,0) circle (0.2);
  \draw (1.3,0) -- (1.7,0);
  \draw (1.5,-0.2) -- (1.5,0.2);
  \node at (0,-2.6) {$\varphi=0$};
  % 红笔：两电荷在中垂线上一点的场强方向及夹角 θ
  \draw[->, red!75!black] (-1.325,0.097) -- (1.2,1.49);
  \draw[->, red!75!black] (1.325,0.097) -- (-1.2,1.49);
  \coordinate (L) at (-1.5,0);
  \coordinate (O) at (0,0);
  \coordinate (Ap) at (0,0.83);
  \pic[draw=red!75!black, angle radius=0.7cm, "\red{$\theta$}", angle eccentricity=1.4] {angle = O--L--Ap};
\end{tikzpicture}
